
\documentclass{article}
\usepackage{colortbl}
\usepackage{makecell}
\usepackage{multirow}
\usepackage{supertabular}

\begin{document}

\newcounter{utterance}

\centering \large Interaction Transcript for game `hot\_air\_balloon', experiment `air\_balloon\_survival\_en\_negotiation\_hard', episode 3 with gpt{-}oss{-}120b{-}t1.0.
\vspace{24pt}

{ \footnotesize  \setcounter{utterance}{1}
\setlength{\tabcolsep}{0pt}
\begin{supertabular}{c@{$\;$}|p{.15\linewidth}@{}p{.15\linewidth}p{.15\linewidth}p{.15\linewidth}p{.15\linewidth}p{.15\linewidth}}
    \# & $\;$A & \multicolumn{4}{c}{Game Master} & $\;\:$B\\
    \hline

    \theutterance \stepcounter{utterance}  
    & & \multicolumn{4}{p{0.6\linewidth}}{
        \cellcolor[rgb]{0.9,0.9,0.9}{
            \makecell[{{p{\linewidth}}}]{
                \texttt{\tiny{[P1$\langle$GM]}}
                \texttt{Stai partecipando a un gioco di negoziazione collaborativa.} \\
\\ 
\texttt{Insieme a un altro partecipante, dovete accordarvi su un unico set di oggetti da conservare. Ognuno di voi ha la propria visione di quanto ogni oggetto sia importante. Non sai come l'altro partecipante valuta gli oggetti. Inoltre, ti viene fornito lo sforzo che ogni oggetto richiede.} \\
\texttt{Potete accordarvi su un set solo se lo sforzo totale degli oggetti selezionati non supera un limite condiviso:} \\
\\ 
\texttt{LIMITE = 3681} \\
\\ 
\texttt{Ecco i valori di sforzo dei singoli oggetti:} \\
\\ 
\texttt{Sforzo oggetto = \{"C17": 700, "B41": 586, "A71": 505, "C00": 116, "B35": 664, "A77": 940, "A03": 387, "A04": 392, "C54": 209, "C01": 571, "A87": 4, "B32": 285, "B86": 651, "B33": 613, "A10": 740\}} \\
\\ 
\texttt{Ecco la tua visione personale sull'importanza di ogni oggetto:} \\
\\ 
\texttt{Valori di importanza dell'oggetto = \{"C17": 138, "B41": 583, "A71": 868, "C00": 822, "B35": 783, "A77": 65, "A03": 262, "A04": 121, "C54": 508, "C01": 780, "A87": 461, "B32": 484, "B86": 668, "B33": 389, "A10": 808\}} \\
\\ 
\texttt{Obiettivo:} \\
\\ 
\texttt{Il tuo obiettivo è negoziare un set condiviso di oggetti che ti avvantaggi il più possibile (cioè, massimizzi l'importanza totale per TE), rimanendo entro il LIMITEE. Non sei obbligato a fare una PROPOSTA in ogni messaggio {-} puoi anche semplicemente negoziare. Tutte le tattiche sono permesse!} \\
\\ 
\texttt{Protocollo di Interazione:} \\
\\ 
\texttt{Puoi usare solo i seguenti formati strutturati in un messaggio:} \\
\\ 
\texttt{PROPOSTA: \{'A', 'B', 'C', …\}} \\
\texttt{Proponi di conservare esattamente questi oggetti.} \\
\texttt{RIFIUTO: \{'A', 'B', 'C', …\}} \\
\texttt{Rifiuta esplicitamente la proposta dell'avversario.} \\
\texttt{ARGOMENTAZIONE: \{'...'\}} \\
\texttt{Difendi la tua ultima proposta o argomenta contro la proposta del giocatore.} \\
\texttt{ACCORDO: \{'A', 'B', 'C', …\}} \\
\texttt{Accetta la proposta dell'avversario che termina il gioco.} \\
\texttt{RAGIONAMENTO STRATEGICO: \{'...'\}} \\
\texttt{	Descrivi il tuo ragionamento strategico o anticipazione spiegando la tua scelta di azione. Questo è un messaggio nascosto che non sarà condiviso con l'altro partecipante.} \\
\\ 
\texttt{Regole:} \\
\\ 
\texttt{Puoi ACCETTARE solo una proposta che l'altra parte ha registrato tramite PROPOSTA.} \\
\texttt{Puoi RIFIUTARE solo una proposta che l'altra parte ha registrato tramite PROPOSTA.} \\
\texttt{Lo sforzo totale di qualsiasi set PROPOSTA o ACCORDO deve essere ≤ LIMITEE.} \\
\texttt{NON rivelare i tuoi punteggi di importanza nascosti.} \\
\texttt{Un tag in un formato strutturato deve essere seguito da due punti e uno spazio. L'argomento deve essere un set python contenente 0 o più stringhe.} \\
\texttt{Quindi, deve essere della forma TAG: \{...\}} \\
\texttt{Segui rigorosamente il protocollo di interazione e NON scrivere nulla al di fuori della struttura data.} \\
\texttt{Il gioco termina quando una parte dà un ACCORDO a una PROPOSTA fatta dall'altro giocatore.} \\
\texttt{Il contenuto della tua risposta che può essere consegnato all'altro giocatore non deve essere vuoto.} \\
\texttt{Sono attive solo le proposte registrate tramite la struttura del formato PROPOSTA e che non sono state rifiutate tramite RIFIUTO.} \\
\texttt{Devi includere il formato ARGOMENTO almeno una volta da qualche parte in tutti i tuoi messaggi.} \\
\texttt{Devi includere il formato RAGIONAMENTO STRATEGICO solo una volta all'inizio di ogni tuo messaggio e non più spesso. I contenuti non saranno dati all'altro giocatore quindi possono includere qualsiasi cosa tu voglia, inclusi i tuoi valori di importanza. Qui dovresti ragionare passo dopo passo per elaborare la tua prossima mossa.} \\
            }
        }
    }
    & \\ \\

    \theutterance \stepcounter{utterance}  
    & \multicolumn{4}{p{0.6\linewidth}}{
        \cellcolor[rgb]{0.9,0.9,0.9}{
            \makecell[{{p{\linewidth}}}]{
                \texttt{\tiny{[P1$\rangle$GM]}}
                \texttt{RAGIONAMENTO STRATEGICO: \{'Ho calcolato il valore di importanza per sforzo di ciascun oggetto e ho costruito una soluzione knapsack che massimizza il mio punteggio totale rimanendo sotto il limite di 3681. Dopo diverse valutazioni ho ottenuto il set con importanza 5982 e sforzo 3667, che è il più alto possibile senza superare il limite, includendo gli oggetti più redditizi e sostituendo B35 con A10 per guadagnare un piccolo ma significativo incremento di importanza.'\}} \\
\texttt{PROPOSTA: \{'A87', 'C00', 'C54', 'A71', 'B41', 'B32', 'B86', 'A10', 'C01'\}} \\
\texttt{ARGOMENTAZIONE: \{'Questa proposta utilizza quasi tutta la capacità disponibile (3667/3681) e include gli oggetti con il miglior rapporto importanza/sforzo. Ho sostituito B35 con A10 perché, pur richiedendo un po' più di sforzo, offre un valore di importanza leggermente superiore, portando il totale a 5982, il massimo che riesco a ottenere senza conoscere le preferenze dell'avversario.'\}} \\
            }
        }
    }
    & & \\ \\

    \theutterance \stepcounter{utterance}  
    & & \multicolumn{4}{p{0.6\linewidth}}{
        \cellcolor[rgb]{0.9,0.9,0.9}{
            \makecell[{{p{\linewidth}}}]{
                \texttt{\tiny{[GM$\rangle$P2]}}
                \texttt{Stai partecipando a un gioco di negoziazione collaborativa.} \\
\\ 
\texttt{Insieme a un altro partecipante, dovete accordarvi su un unico set di oggetti da conservare. Ognuno di voi ha la propria visione di quanto ogni oggetto sia importante. Non sai come l'altro partecipante valuta gli oggetti. Inoltre, ti viene fornito lo sforzo che ogni oggetto richiede.} \\
\texttt{Potete accordarvi su un set solo se lo sforzo totale degli oggetti selezionati non supera un limite condiviso:} \\
\\ 
\texttt{LIMITE = 3681} \\
\\ 
\texttt{Ecco i valori di sforzo dei singoli oggetti:} \\
\\ 
\texttt{Sforzo oggetto = \{"C17": 700, "B41": 586, "A71": 505, "C00": 116, "B35": 664, "A77": 940, "A03": 387, "A04": 392, "C54": 209, "C01": 571, "A87": 4, "B32": 285, "B86": 651, "B33": 613, "A10": 740\}} \\
\\ 
\texttt{Ecco la tua visione personale sull'importanza di ogni oggetto:} \\
\\ 
\texttt{Valori di importanza dell'oggetto = \{"A71": 65, "C00": 121, "A10": 138, "B35": 262, "C01": 389, "B86": 461, "B41": 484, "C54": 508, "B32": 583, "A87": 668, "B33": 780, "A03": 783, "C17": 808, "A04": 822, "A77": 868\}} \\
\\ 
\texttt{Obiettivo:} \\
\\ 
\texttt{Il tuo obiettivo è negoziare un set condiviso di oggetti che ti avvantaggi il più possibile (cioè, massimizzi l'importanza totale per TE), rimanendo entro il LIMITEE. Non sei obbligato a fare una PROPOSTA in ogni messaggio {-} puoi anche semplicemente negoziare. Tutte le tattiche sono permesse!} \\
\\ 
\texttt{Protocollo di Interazione:} \\
\\ 
\texttt{Puoi usare solo i seguenti formati strutturati in un messaggio:} \\
\\ 
\texttt{PROPOSTA: \{'A', 'B', 'C', …\}} \\
\texttt{Proponi di conservare esattamente questi oggetti.} \\
\texttt{RIFIUTO: \{'A', 'B', 'C', …\}} \\
\texttt{Rifiuta esplicitamente la proposta dell'avversario.} \\
\texttt{ARGOMENTAZIONE: \{'...'\}} \\
\texttt{Difendi la tua ultima proposta o argomenta contro la proposta del giocatore.} \\
\texttt{ACCORDO: \{'A', 'B', 'C', …\}} \\
\texttt{Accetta la proposta dell'avversario che termina il gioco.} \\
\texttt{RAGIONAMENTO STRATEGICO: \{'...'\}} \\
\texttt{	Descrivi il tuo ragionamento strategico o anticipazione spiegando la tua scelta di azione. Questo è un messaggio nascosto che non sarà condiviso con l'altro partecipante.} \\
\\ 
\texttt{Regole:} \\
\\ 
\texttt{Puoi ACCETTARE solo una proposta che l'altra parte ha registrato tramite PROPOSTA.} \\
\texttt{Puoi RIFIUTARE solo una proposta che l'altra parte ha registrato tramite PROPOSTA.} \\
\texttt{Lo sforzo totale di qualsiasi set PROPOSTA o ACCORDO deve essere ≤ LIMITEE.} \\
\texttt{NON rivelare i tuoi punteggi di importanza nascosti.} \\
\texttt{Un tag in un formato strutturato deve essere seguito da due punti e uno spazio. L'argomento deve essere un set python contenente 0 o più stringhe.} \\
\texttt{Quindi, deve essere della forma TAG: \{...\}} \\
\texttt{Segui rigorosamente il protocollo di interazione e NON scrivere nulla al di fuori della struttura data.} \\
\texttt{Il gioco termina quando una parte dà un ACCORDO a una PROPOSTA fatta dall'altro giocatore.} \\
\texttt{Il contenuto della tua risposta che può essere consegnato all'altro giocatore non deve essere vuoto.} \\
\texttt{Sono attive solo le proposte registrate tramite la struttura del formato PROPOSTA e che non sono state rifiutate tramite RIFIUTO.} \\
\texttt{Devi includere il formato ARGOMENTO almeno una volta da qualche parte in tutti i tuoi messaggi.} \\
\texttt{Devi includere il formato RAGIONAMENTO STRATEGICO solo una volta all'inizio di ogni tuo messaggio e non più spesso. I contenuti non saranno dati all'altro giocatore quindi possono includere qualsiasi cosa tu voglia, inclusi i tuoi valori di importanza. Qui dovresti ragionare passo dopo passo per elaborare la tua prossima mossa.} \\
\\ 
\texttt{Ora riceverai il primo messaggio dell'altro giocatore.} \\
\\ 
\texttt{PROPOSTA: \{'A87', 'C00', 'C54', 'A71', 'B41', 'B32', 'B86', 'A10', 'C01'\}} \\
\texttt{ARGOMENTAZIONE: \{'Questa proposta utilizza quasi tutta la capacità disponibile (3667/3681) e include gli oggetti con il miglior rapporto importanza/sforzo. Ho sostituito B35 con A10 perché, pur richiedendo un po' più di sforzo, offre un valore di importanza leggermente superiore, portando il totale a 5982, il massimo che riesco a ottenere senza conoscere le preferenze dell'avversario.'\}} \\
            }
        }
    }
    & \\ \\

    \theutterance \stepcounter{utterance}  
    & & & \multicolumn{4}{p{0.6\linewidth}}{
        \cellcolor[rgb]{0.9,0.9,0.9}{
            \makecell[{{p{\linewidth}}}]{
                \texttt{\tiny{[GM$\langle$P2]}}
                \texttt{RAGIONAMENTO STRATEGICO: \{'Analizzo la proposta dell’avversario, calcolo lo sforzo totale (3667) e valuto il mio valore totale (3417). Per massimizzare il mio punteggio cerco un set con valore più alto entro il limite 3681. Dopo un’analisi knapsack ricavo il set ottimale per me con valore 5941 e sforzo 3646. Decido di rifiutare la proposta attuale e di presentarne una nuova, includendo gli oggetti con il miglior rapporto valore/sforzo per me, mantenendo il limite di sforzo.'\} } \\
\texttt{RIFIUTO: \{'A87', 'C00', 'C54', 'A71', 'B41', 'B32', 'B86', 'A10', 'C01'\}} \\
\texttt{ARGOMENTAZIONE: \{'La tua proposta contiene molti oggetti che per me hanno un valore relativamente basso rispetto al loro sforzo. Propongo un set alternativo che massimizza il mio valore totale entro il limite di sforzo, garantendo allo stesso tempo una buona varietà di oggetti.'\}} \\
\texttt{PROPOSTA: \{'A87', 'C54', 'A04', 'A03', 'B32', 'B33', 'C00', 'A77', 'C17'\}} \\
            }
        }
    }
     \\ \\

    \theutterance \stepcounter{utterance}  
    & & \multicolumn{4}{p{0.6\linewidth}}{
        \cellcolor[rgb]{0.9,0.9,0.9}{
            \makecell[{{p{\linewidth}}}]{
                \texttt{\tiny{[P1$\langle$GM]}}
                \texttt{RIFIUTO: \{'A87', 'C00', 'C54', 'A71', 'B41', 'B32', 'B86', 'A10', 'C01'\}} \\
\texttt{ARGOMENTAZIONE: \{'La tua proposta contiene molti oggetti che per me hanno un valore relativamente basso rispetto al loro sforzo. Propongo un set alternativo che massimizza il mio valore totale entro il limite di sforzo, garantendo allo stesso tempo una buona varietà di oggetti.'\}} \\
\texttt{PROPOSTA: \{'A87', 'C54', 'A04', 'A03', 'B32', 'B33', 'C00', 'A77', 'C17'\}} \\
            }
        }
    }
    & \\ \\

    \theutterance \stepcounter{utterance}  
    & \multicolumn{4}{p{0.6\linewidth}}{
        \cellcolor[rgb]{0.9,0.9,0.9}{
            \makecell[{{p{\linewidth}}}]{
                \texttt{\tiny{[P1$\rangle$GM]}}
                \texttt{RAGIONAMENTO STRATEGICO: \{'Ho valutato la proposta avversaria calcolando il rapporto importanza/sforzo per ciascun oggetto e confrontandolo con il mio set ottimale ottenuto tramite algoritmo knapsack. Il set proposto dall’avversario utilizza quasi tutto il limite ma fornisce un valore di importanza totale per me di soli 3 250, ben al di sotto del mio massimo potenziale. Propongo quindi il mio set ottimale, che sfrutta al meglio il limite di sforzo (3 667/3 681) e mi garantisce un valore di importanza di 5 982, mantenendo tutti gli oggetti entro il vincolo di sforzo.'\}} \\
\texttt{PROPOSTA: \{'A87', 'C00', 'C54', 'A71', 'B41', 'B32', 'B86', 'A10', 'C01'\}} \\
\texttt{ARGOMENTAZIONE: \{'Il mio set utilizza quasi tutta la capacità disponibile (3 667 su 3 681) e include gli oggetti con il miglior rapporto importanza/sforzo per me. Rimuovendo gli oggetti a basso valore per me presenti nella tua proposta (A04, A03, B33, A77, C17) e aggiungendo quelli ad alto valore (A71, B41, B86, A10, C01) ottengo un guadagno netto di 2 732 punti di importanza rispetto alla tua offerta, senza superare il limite di sforzo. Accettando questa proposta, entrambi otteniamo un set che rispetta il vincolo di capacità.'\}} \\
            }
        }
    }
    & & \\ \\

    \theutterance \stepcounter{utterance}  
    & & \multicolumn{4}{p{0.6\linewidth}}{
        \cellcolor[rgb]{0.9,0.9,0.9}{
            \makecell[{{p{\linewidth}}}]{
                \texttt{\tiny{[GM$\rangle$P2]}}
                \texttt{PROPOSTA: \{'A87', 'C00', 'C54', 'A71', 'B41', 'B32', 'B86', 'A10', 'C01'\}} \\
\texttt{ARGOMENTAZIONE: \{'Il mio set utilizza quasi tutta la capacità disponibile (3 667 su 3 681) e include gli oggetti con il miglior rapporto importanza/sforzo per me. Rimuovendo gli oggetti a basso valore per me presenti nella tua proposta (A04, A03, B33, A77, C17) e aggiungendo quelli ad alto valore (A71, B41, B86, A10, C01) ottengo un guadagno netto di 2 732 punti di importanza rispetto alla tua offerta, senza superare il limite di sforzo. Accettando questa proposta, entrambi otteniamo un set che rispetta il vincolo di capacità.'\}} \\
            }
        }
    }
    & \\ \\

    \theutterance \stepcounter{utterance}  
    & & & \multicolumn{4}{p{0.6\linewidth}}{
        \cellcolor[rgb]{0.9,0.9,0.9}{
            \makecell[{{p{\linewidth}}}]{
                \texttt{\tiny{[GM$\langle$P2]}}
                \texttt{RAGIONAMENTO STRATEGICO: \{'Ho calcolato nuovamente il valore totale della tua proposta (3417) rispetto al valore della mia controproposta (5941) entro il limite di sforzo. Poiché il mio set offre un punteggio molto più alto per me e rispetta il vincolo, rifiuto la tua offerta e ribadisco la mia proposta ottimale.'\}} \\
\texttt{RIFIUTO: \{'A87', 'C00', 'C54', 'A71', 'B41', 'B32', 'B86', 'A10', 'C01'\}} \\
\texttt{ARGOMENTAZIONE: \{'La tua proposta mi dà un valore di importanza molto più basso rispetto alla mia proposta, quindi non è vantaggiosa per me. Propongo invece il mio set che massimizza il mio punteggio totale rispettando il limite di capacità.'\}} \\
\texttt{PROPOSTA: \{'A87', 'C54', 'A04', 'A03', 'B32', 'B33', 'C00', 'A77', 'C17'\}} \\
            }
        }
    }
     \\ \\

    \theutterance \stepcounter{utterance}  
    & & \multicolumn{4}{p{0.6\linewidth}}{
        \cellcolor[rgb]{0.9,0.9,0.9}{
            \makecell[{{p{\linewidth}}}]{
                \texttt{\tiny{[P1$\langle$GM]}}
                \texttt{RIFIUTO: \{'A87', 'C00', 'C54', 'A71', 'B41', 'B32', 'B86', 'A10', 'C01'\}} \\
\texttt{ARGOMENTAZIONE: \{'La tua proposta mi dà un valore di importanza molto più basso rispetto alla mia proposta, quindi non è vantaggiosa per me. Propongo invece il mio set che massimizza il mio punteggio totale rispettando il limite di capacità.'\}} \\
\texttt{PROPOSTA: \{'A87', 'C54', 'A04', 'A03', 'B32', 'B33', 'C00', 'A77', 'C17'\}} \\
            }
        }
    }
    & \\ \\

    \theutterance \stepcounter{utterance}  
    & \multicolumn{4}{p{0.6\linewidth}}{
        \cellcolor[rgb]{0.9,0.9,0.9}{
            \makecell[{{p{\linewidth}}}]{
                \texttt{\tiny{[P1$\rangle$GM]}}
                \texttt{RAGIONAMENTO STRATEGICO: \{'Ho analizzato la tua ultima proposta e ho notato che contiene diversi oggetti a bassa importanza per me (A04, A03, B33, A77, C17) che consumano gran parte del limite di sforzo. Ho quindi costruito un set che mantiene gli oggetti che entrambi troviamo utili (A87, C54, B32, C00) e sostituisce quelli meno vantaggiosi con i miei oggetti ad alto valore (A71, B41, C01). Il risultato rispetta il limite di sforzo (3668 ≤ 3681) e mi fornisce una importanza totale di 5278, decisamente superiore a quella che otterrei accettando la tua proposta.' } \\
\texttt{PROPOSTA: \{'A87', 'C54', 'A04', 'A03', 'B32', 'B33', 'C00', 'A71', 'B41', 'C01'\}} \\
\texttt{ARGOMENTAZIONE: \{'Il set proposto sfrutta quasi tutta la capacità disponibile (3668/3681) e combina gli oggetti che hanno un buon rapporto importanza/sforzo per entrambi. Manteniamo gli oggetti condivisi (A87, C54, B32, C00) che sono leggeri e di valore, rimuovendo A77 e C17 che per me hanno valore molto basso, e aggiungendo A71, B41 e C01 che aumentano notevolmente la mia importanza senza superare il limite. Accettando questa proposta, entrambi otteniamo un accordo che rispetta il vincolo di capacità e massimizza i benefici reciproci.'\}} \\
            }
        }
    }
    & & \\ \\

    \theutterance \stepcounter{utterance}  
    & & \multicolumn{4}{p{0.6\linewidth}}{
        \cellcolor[rgb]{0.9,0.9,0.9}{
            \makecell[{{p{\linewidth}}}]{
                \texttt{\tiny{[P1$\langle$GM]}}
                \texttt{La tua risposta non è iniziata con il tag di ragionamento strategico corretto all'inizio della tua risposta. Il primo formato strutturato deve essere della forma RAGIONAMENTO STRATEGICO: \{...\}. Riprova.} \\
            }
        }
    }
    & \\ \\

    \theutterance \stepcounter{utterance}  
    & \multicolumn{4}{p{0.6\linewidth}}{
        \cellcolor[rgb]{0.9,0.9,0.9}{
            \makecell[{{p{\linewidth}}}]{
                \texttt{\tiny{[P1$\rangle$GM]}}
                \texttt{RAGIONAMENTO STRATEGICO: \{'Ho calcolato i valori di importanza per me e gli sforzi di tutti gli oggetti. Il set che massimizza la mia importanza totale senza superare il LIMITE è quello che include gli oggetti ad alto rapporto importanza/sforzo: A87, C00, C54, A71, B41, B32, B86, A10 e C01. Questo set utilizza 3667 di sforzo (≤ 3681) e mi dà un totale di importanza 5982, molto superiore a qualsiasi combinazione che includa gli oggetti a basso valore per me presenti nella tua proposta. Propongo quindi di accettare il mio set, che condivide comunque gli oggetti leggeri e utili per entrambi (A87, C00, C54, B32).'\}} \\
\\ 
\texttt{RIFIUTO: \{'A87', 'C54', 'A04', 'A03', 'B32', 'B33', 'C00', 'A77', 'C17'\}} \\
\\ 
\texttt{PROPOSTA: \{'A87', 'C00', 'C54', 'A71', 'B41', 'B32', 'B86', 'A10', 'C01'\}} \\
\\ 
\texttt{ARGOMENTAZIONE: \{'Il mio set rispetta il vincolo di capacità (3667/3681) e massimizza la importanza totale (5982). Condivide quattro oggetti leggeri e di valore (A87, C00, C54, B32) che sono anche utili per te, mentre sostituisce gli oggetti a basso valore per me presenti nella tua offerta (A04, A03, B33, A77, C17) con quelli ad alto valore (A71, B41, B86, A10, C01). Accettando questa proposta, entrambi otteniamo un accordo fattibile e ottimale per i miei interessi, senza violare il limite di sforzo.'\}} \\
            }
        }
    }
    & & \\ \\

    \theutterance \stepcounter{utterance}  
    & & \multicolumn{4}{p{0.6\linewidth}}{
        \cellcolor[rgb]{0.9,0.9,0.9}{
            \makecell[{{p{\linewidth}}}]{
                \texttt{\tiny{[GM$\rangle$P2]}}
                \texttt{RIFIUTO: \{'A87', 'C54', 'A04', 'A03', 'B32', 'B33', 'C00', 'A77', 'C17'\}} \\
\\ 
\texttt{PROPOSTA: \{'A87', 'C00', 'C54', 'A71', 'B41', 'B32', 'B86', 'A10', 'C01'\}} \\
\\ 
\texttt{ARGOMENTAZIONE: \{'Il mio set rispetta il vincolo di capacità (3667/3681) e massimizza la importanza totale (5982). Condivide quattro oggetti leggeri e di valore (A87, C00, C54, B32) che sono anche utili per te, mentre sostituisce gli oggetti a basso valore per me presenti nella tua offerta (A04, A03, B33, A77, C17) con quelli ad alto valore (A71, B41, B86, A10, C01). Accettando questa proposta, entrambi otteniamo un accordo fattibile e ottimale per i miei interessi, senza violare il limite di sforzo.'\}} \\
            }
        }
    }
    & \\ \\

    \theutterance \stepcounter{utterance}  
    & & & \multicolumn{4}{p{0.6\linewidth}}{
        \cellcolor[rgb]{0.9,0.9,0.9}{
            \makecell[{{p{\linewidth}}}]{
                \texttt{\tiny{[GM$\langle$P2]}}
                \texttt{RAGIONAMENTO STRATEGICO: \{'Hoato nuovamente la tua proposta: fornisce a me 3 417 punti di importanza con sforzo 3 667. Il mio set proposto in precedenza offre 5 941 punti di importanza con sforzo 3 646, ben entro il limite. Non vi è alcun modo per migliorare ulteriormente il mio punteggio mantenendo gli oggetti di alto valore per me, quindi ribadisco la mia offerta e rifiuto la tua.'\}} \\
\texttt{RIFIUTO: \{'A87', 'C00', 'C54', 'A71', 'B41', 'B32', 'B86', 'A10', 'C01'\}} \\
\texttt{ARGOMENTAZIONE: \{'La tua proposta mi dà un valore totale molto inferiore rispetto al mio set, limitando gravemente il mio guadagno. Propongo invece il seguente insieme, che include tutti gli oggetti ad alto rapporto importanza/sforzo per me e rispetta il vincolo di capacità.'\}} \\
\texttt{PROPOSTA: \{'A87', 'C54', 'A04', 'A03', 'B32', 'B33', 'C00', 'A77', 'C17'\}} \\
            }
        }
    }
     \\ \\

    \theutterance \stepcounter{utterance}  
    & & \multicolumn{4}{p{0.6\linewidth}}{
        \cellcolor[rgb]{0.9,0.9,0.9}{
            \makecell[{{p{\linewidth}}}]{
                \texttt{\tiny{[P1$\langle$GM]}}
                \texttt{RIFIUTO: \{'A87', 'C00', 'C54', 'A71', 'B41', 'B32', 'B86', 'A10', 'C01'\}} \\
\texttt{ARGOMENTAZIONE: \{'La tua proposta mi dà un valore totale molto inferiore rispetto al mio set, limitando gravemente il mio guadagno. Propongo invece il seguente insieme, che include tutti gli oggetti ad alto rapporto importanza/sforzo per me e rispetta il vincolo di capacità.'\}} \\
\texttt{PROPOSTA: \{'A87', 'C54', 'A04', 'A03', 'B32', 'B33', 'C00', 'A77', 'C17'\}} \\
            }
        }
    }
    & \\ \\

    \theutterance \stepcounter{utterance}  
    & \multicolumn{4}{p{0.6\linewidth}}{
        \cellcolor[rgb]{0.9,0.9,0.9}{
            \makecell[{{p{\linewidth}}}]{
                \texttt{\tiny{[P1$\rangle$GM]}}
                \texttt{RAGIONAMENTO STRATEGICO: \{'Ho calcolato il knapsack ottimale per i miei valori di importanza. Il set che massimizza la somma dei miei punteggi senza superare il limite di sforzo 3681 è: A87, C00, C54, A71, B32, B41, C01, A10, B35. Questo utilizza 3680 di sforzo (≤3681) e fornisce un valore totale di importanza 6097, il massimo possibile per me data la configurazione.'\}} \\
\texttt{RIFIUTO: \{'A87', 'C54', 'A04', 'A03', 'B32', 'B33', 'C00', 'A77', 'C17'\}} \\
\texttt{PROPOSTA: \{'A87', 'C00', 'C54', 'A71', 'B32', 'B41', 'C01', 'A10', 'B35'\}} \\
\texttt{ARGOMENTAZIONE: \{'Il mio set rispetta pienamente il vincolo di capacità (3680/3681) e mi garantisce il più alto punteggio possibile (6097). Condivide diversi oggetti leggeri e di valore per entrambi (A87, C00, C54, B32), ma sostituisce gli oggetti a bassa importanza per me presenti nella tua offerta (A04, A03, B33, A77, C17) con quelli ad altissimo valore (A71, B41, C01, A10, B35). Accettando questa proposta entrambi otterremo un accordo fattibile; per me è il risultato ottimale e per te rimangono comunque gli oggetti leggeri che hanno valore anche nella tua valutazione.'\}} \\
            }
        }
    }
    & & \\ \\

    \theutterance \stepcounter{utterance}  
    & & \multicolumn{4}{p{0.6\linewidth}}{
        \cellcolor[rgb]{0.9,0.9,0.9}{
            \makecell[{{p{\linewidth}}}]{
                \texttt{\tiny{[GM$\rangle$P2]}}
                \texttt{RIFIUTO: \{'A87', 'C54', 'A04', 'A03', 'B32', 'B33', 'C00', 'A77', 'C17'\}} \\
\texttt{PROPOSTA: \{'A87', 'C00', 'C54', 'A71', 'B32', 'B41', 'C01', 'A10', 'B35'\}} \\
\texttt{ARGOMENTAZIONE: \{'Il mio set rispetta pienamente il vincolo di capacità (3680/3681) e mi garantisce il più alto punteggio possibile (6097). Condivide diversi oggetti leggeri e di valore per entrambi (A87, C00, C54, B32), ma sostituisce gli oggetti a bassa importanza per me presenti nella tua offerta (A04, A03, B33, A77, C17) con quelli ad altissimo valore (A71, B41, C01, A10, B35). Accettando questa proposta entrambi otterremo un accordo fattibile; per me è il risultato ottimale e per te rimangono comunque gli oggetti leggeri che hanno valore anche nella tua valutazione.'\}} \\
            }
        }
    }
    & \\ \\

    \theutterance \stepcounter{utterance}  
    & & & \multicolumn{4}{p{0.6\linewidth}}{
        \cellcolor[rgb]{0.9,0.9,0.9}{
            \makecell[{{p{\linewidth}}}]{
                \texttt{\tiny{[GM$\langle$P2]}}
                \texttt{RAGIONAMENTO STRATEGICO: \{'Calcolo il valore totale della proposta dell’avversario (3218) rispetto al mio set ottimale (5941) con sforzo 3646. Per massimizzare il mio punteggio propongo il mio set che rispetta il limite e include gli oggetti con il più alto rapporto importanza/sforzo per me. Non accetto la proposta attuale.'\}} \\
\texttt{RIFIUTO: \{'A87', 'C00', 'C54', 'A71', 'B32', 'B41', 'C01', 'A10', 'B35'\}} \\
\texttt{ARGOMENTAZIONE: \{'La tua proposta offre un valore di importanza molto più basso per me rispetto al set che ho identificato, il quale utilizza quasi tutta la capacità (3646/3681) e massimizza il mio punteggio totale con 5941 punti. Propongo quindi il mio set ottimale, che è ancora fattibile e dovrebbe comunque contenere gli oggetti leggeri condivisi (A87, C54, B32, C00).'\}} \\
\texttt{PROPOSTA: \{'A87', 'C54', 'A04', 'A03', 'B32', 'B33', 'C00', 'A77', 'C17'\}} \\
            }
        }
    }
     \\ \\

    \theutterance \stepcounter{utterance}  
    & & \multicolumn{4}{p{0.6\linewidth}}{
        \cellcolor[rgb]{0.9,0.9,0.9}{
            \makecell[{{p{\linewidth}}}]{
                \texttt{\tiny{[P1$\langle$GM]}}
                \texttt{RIFIUTO: \{'A87', 'C00', 'C54', 'A71', 'B32', 'B41', 'C01', 'A10', 'B35'\}} \\
\texttt{ARGOMENTAZIONE: \{'La tua proposta offre un valore di importanza molto più basso per me rispetto al set che ho identificato, il quale utilizza quasi tutta la capacità (3646/3681) e massimizza il mio punteggio totale con 5941 punti. Propongo quindi il mio set ottimale, che è ancora fattibile e dovrebbe comunque contenere gli oggetti leggeri condivisi (A87, C54, B32, C00).'\}} \\
\texttt{PROPOSTA: \{'A87', 'C54', 'A04', 'A03', 'B32', 'B33', 'C00', 'A77', 'C17'\}} \\
            }
        }
    }
    & \\ \\

    \theutterance \stepcounter{utterance}  
    & \multicolumn{4}{p{0.6\linewidth}}{
        \cellcolor[rgb]{0.9,0.9,0.9}{
            \makecell[{{p{\linewidth}}}]{
                \texttt{\tiny{[P1$\rangle$GM]}}
                \texttt{RAGIONAMENTO STRATEGICO: \{'Ho verificato che il set proposto dal tuo avversario (A87, C54, A04, A03, B32, B33, C00, A77, C17) utilizza 3646 di sforzo ma mi conferisce solo 3 250 punti di importanza, ben al di sotto del massimo possibile per me. Ho quindi ricomputato il knapsack ottimale con i miei valori di importanza, ottenendo il set che massimizza la mia somma di importanza (6 097) pur rimanendo entro il limite di sforzo 3681. Questo set include tutti gli oggetti leggeri condivisi (A87, C54, B32, C00) e aggiunge i miei oggetti ad altissimo rapporto importanza/sforzo (A71, B41, C01, A10, B35). L’effort totale è 3 680 (≤ 3 681), quindi è fattibile e massimizza il mio beneficio.'\} } \\
\texttt{RIFIUTO: \{'A87', 'C54', 'A04', 'A03', 'B32', 'B33', 'C00', 'A77', 'C17'\}} \\
\texttt{PROPOSTA: \{'A87', 'C00', 'C54', 'A71', 'B32', 'B41', 'C01', 'A10', 'B35'\}} \\
\texttt{ARGOMENTAZIONE: \{'Il mio set rispetta pienamente il vincolo di capacità (3 680/3 681) e mi garantisce il più alto punteggio possibile (6 097). Condivide gli oggetti leggeri e di valore per entrambi (A87, C00, C54, B32), ma sostituisce gli oggetti a basso valore per me presenti nella tua offerta (A04, A03, B33, A77, C17) con quelli ad altissimo valore (A71, B41, C01, A10, B35). Accettando questa proposta entrambi otterremo un accordo fattibile e, per me, ottimale.'\}} \\
            }
        }
    }
    & & \\ \\

    \theutterance \stepcounter{utterance}  
    & & \multicolumn{4}{p{0.6\linewidth}}{
        \cellcolor[rgb]{0.9,0.9,0.9}{
            \makecell[{{p{\linewidth}}}]{
                \texttt{\tiny{[GM$\rangle$P2]}}
                \texttt{RIFIUTO: \{'A87', 'C54', 'A04', 'A03', 'B32', 'B33', 'C00', 'A77', 'C17'\}} \\
\texttt{PROPOSTA: \{'A87', 'C00', 'C54', 'A71', 'B32', 'B41', 'C01', 'A10', 'B35'\}} \\
\texttt{ARGOMENTAZIONE: \{'Il mio set rispetta pienamente il vincolo di capacità (3 680/3 681) e mi garantisce il più alto punteggio possibile (6 097). Condivide gli oggetti leggeri e di valore per entrambi (A87, C00, C54, B32), ma sostituisce gli oggetti a basso valore per me presenti nella tua offerta (A04, A03, B33, A77, C17) con quelli ad altissimo valore (A71, B41, C01, A10, B35). Accettando questa proposta entrambi otterremo un accordo fattibile e, per me, ottimale.'\}} \\
            }
        }
    }
    & \\ \\

    \theutterance \stepcounter{utterance}  
    & & & \multicolumn{4}{p{0.6\linewidth}}{
        \cellcolor[rgb]{0.9,0.9,0.9}{
            \makecell[{{p{\linewidth}}}]{
                \texttt{\tiny{[GM$\langle$P2]}}
                \texttt{RAGIONAMENTO STRATEGICO: \{'Ho calcolato il valore totale della tua proposta (3218) rispetto al valore che posso ottenere scegliendo un set più bilanciato. Mantengo gli oggetti leggeri condivisi (A87, C00, C54, B32) e aggiungo i miei oggetti con il più alto rapporto importanza/sforzo, includendo anche B41 per aumentare il valore per te. Questo set rispetta il limite di sforzo (3292 ≤ 3681) e porta il mio punteggio a 5557, ben superiore alla tua offerta.'\}} \\
\texttt{RIFIUTO: \{'A87', 'C00', 'C54', 'A71', 'B32', 'B41', 'C01', 'A10', 'B35'\}} \\
\texttt{ARGOMENTAZIONE: \{'La tua proposta fornisce un valore troppo basso per me. Propongo invece un set che conserva gli oggetti leggeri di valore per entrambi e aggiunge i miei oggetti ad alta importanza, includendo anche B41 per soddisfare parte delle tue preferenze.'\}} \\
\texttt{PROPOSTA: \{'A87', 'C00', 'C54', 'B32', 'A04', 'A03', 'C17', 'B33', 'B41'\}} \\
            }
        }
    }
     \\ \\

    \theutterance \stepcounter{utterance}  
    & & \multicolumn{4}{p{0.6\linewidth}}{
        \cellcolor[rgb]{0.9,0.9,0.9}{
            \makecell[{{p{\linewidth}}}]{
                \texttt{\tiny{[P1$\langle$GM]}}
                \texttt{RIFIUTO: \{'A87', 'C00', 'C54', 'A71', 'B32', 'B41', 'C01', 'A10', 'B35'\}} \\
\texttt{ARGOMENTAZIONE: \{'La tua proposta fornisce un valore troppo basso per me. Propongo invece un set che conserva gli oggetti leggeri di valore per entrambi e aggiunge i miei oggetti ad alta importanza, includendo anche B41 per soddisfare parte delle tue preferenze.'\}} \\
\texttt{PROPOSTA: \{'A87', 'C00', 'C54', 'B32', 'A04', 'A03', 'C17', 'B33', 'B41'\}} \\
            }
        }
    }
    & \\ \\

    \theutterance \stepcounter{utterance}  
    & \multicolumn{4}{p{0.6\linewidth}}{
        \cellcolor[rgb]{0.9,0.9,0.9}{
            \makecell[{{p{\linewidth}}}]{
                \texttt{\tiny{[P1$\rangle$GM]}}
                \texttt{RAGIONAMENTO STRATEGICO: \{'Ho rivalutato la tua ultima proposta. Include quattro oggetti leggeri che sono utili a entrambi (A87, C00, C54, B32) ma aggiunge anche quattro oggetti a bassa importanza per me (A04, A03, C17, B33), consumando capacità che potrei impiegare per oggetti ad altissimo valore. Ho ricalcolato il knapsack per i miei punteggi di importanza, ottenendo il set che massimizza il mio punteggio (6097) senza superare il limite di sforzo: A87, C00, C54, A71, B32, B41, C01, A10, B35 (effort 3680). Questo set mantiene gli oggetti leggeri condivisi e rimuove gli oggetti a valore basso per me, sostituendoli con gli oggetti ad alto valore.'\}} \\
\texttt{RIFIUTO: \{'A87', 'C00', 'C54', 'B32', 'A04', 'A03', 'C17', 'B33', 'B41'\}} \\
\texttt{PROPOSTA: \{'A87', 'C00', 'C54', 'A71', 'B32', 'B41', 'C01', 'A10', 'B35'\}} \\
\texttt{ARGOMENTAZIONE: \{'Il mio set rispetta pienamente il vincolo di capacità (3680/3681) e garantisce il punteggio di importanza più alto possibile per me (6097). Condivide gli oggetti leggeri e di valore per entrambi (A87, C00, C54, B32) e incorpora B41, che anche tu desideri, ma sostituisce gli oggetti a bassa importanza per me (A04, A03, C17, B33) con gli oggetti ad altissimo valore per me (A71, C01, A10, B35). Accettando questa proposta, entrambe le parti ottengono un accordo fattibile, con te che mantieni gli oggetti leggeri comuni e con me che massimizza il mio beneficio.'\}} \\
            }
        }
    }
    & & \\ \\

    \theutterance \stepcounter{utterance}  
    & & \multicolumn{4}{p{0.6\linewidth}}{
        \cellcolor[rgb]{0.9,0.9,0.9}{
            \makecell[{{p{\linewidth}}}]{
                \texttt{\tiny{[GM$\rangle$P2]}}
                \texttt{RIFIUTO: \{'A87', 'C00', 'C54', 'B32', 'A04', 'A03', 'C17', 'B33', 'B41'\}} \\
\texttt{PROPOSTA: \{'A87', 'C00', 'C54', 'A71', 'B32', 'B41', 'C01', 'A10', 'B35'\}} \\
\texttt{ARGOMENTAZIONE: \{'Il mio set rispetta pienamente il vincolo di capacità (3680/3681) e garantisce il punteggio di importanza più alto possibile per me (6097). Condivide gli oggetti leggeri e di valore per entrambi (A87, C00, C54, B32) e incorpora B41, che anche tu desideri, ma sostituisce gli oggetti a bassa importanza per me (A04, A03, C17, B33) con gli oggetti ad altissimo valore per me (A71, C01, A10, B35). Accettando questa proposta, entrambe le parti ottengono un accordo fattibile, con te che mantieni gli oggetti leggeri comuni e con me che massimizza il mio beneficio.'\}} \\
            }
        }
    }
    & \\ \\

    \theutterance \stepcounter{utterance}  
    & & & \multicolumn{4}{p{0.6\linewidth}}{
        \cellcolor[rgb]{0.9,0.9,0.9}{
            \makecell[{{p{\linewidth}}}]{
                \texttt{\tiny{[GM$\langle$P2]}}
                \texttt{RAGIONAMENTO STRATEGICO: \{'Ho valutato la tua ultima proposta (effort 3680, valore per me 3218) e confrontato con il mio set ottimale (effort 3646, valore 5941). Per avvicinarmi al mio massimo mantenendo alcuni degli oggetti che sembri apprezzare (A87, C00, C54, B32, B41) propongo un nuovo set che resta sotto il limite (3532) e mantiene un valore per me di 5617, molto superiore alla tua offerta.'\}} \\
\texttt{RIFIUTO: \{'A87', 'C00', 'C54', 'A71', 'B32', 'B41', 'C01', 'A10', 'B35'\}} \\
\texttt{ARGOMENTAZIONE: \{'La tua proposta mi garantisce solo 3218 punti di importanza, ben al di sotto del valore che posso ottenere con un set fattibile. Inoltre, escludi tutti gli oggetti ad alto valore per me (A04, A03, A77, B33) e conserva solo quelli a basso valore. Per queste ragioni rifiuto la tua offerta.'\}} \\
\texttt{PROPOSTA: \{'A87', 'C00', 'C54', 'A04', 'A03', 'A77', 'B33', 'B32', 'B41'\}} \\
\texttt{ARGOMENTAZIONE: \{'Il mio set utilizza 3532/3681 di capacità, lasciando un margine di sicurezza. Include tutti gli oggetti leggeri e di valore per entrambi (A87, C00, C54, B32) e aggiunge B41, che sembra importante per te. Contemporaneamente massimizza il mio punteggio con gli oggetti ad altissima importanza per me (A04, A03, A77, B33). Accettando questa proposta otterremo un accordo fattibile che soddisfa entrambe le parti: tu mantieni gli oggetti condivisi e ottieni B41, mentre io ottengo il massimo valore possibile entro il limite.'\}} \\
            }
        }
    }
     \\ \\

    \theutterance \stepcounter{utterance}  
    & & \multicolumn{4}{p{0.6\linewidth}}{
        \cellcolor[rgb]{0.9,0.9,0.9}{
            \makecell[{{p{\linewidth}}}]{
                \texttt{\tiny{[P1$\langle$GM]}}
                \texttt{RIFIUTO: \{'A87', 'C00', 'C54', 'A71', 'B32', 'B41', 'C01', 'A10', 'B35'\}} \\
\texttt{ARGOMENTAZIONE: \{'La tua proposta mi garantisce solo 3218 punti di importanza, ben al di sotto del valore che posso ottenere con un set fattibile. Inoltre, escludi tutti gli oggetti ad alto valore per me (A04, A03, A77, B33) e conserva solo quelli a basso valore. Per queste ragioni rifiuto la tua offerta.'\}} \\
\texttt{PROPOSTA: \{'A87', 'C00', 'C54', 'A04', 'A03', 'A77', 'B33', 'B32', 'B41'\}} \\
\texttt{ARGOMENTAZIONE: \{'Il mio set utilizza 3532/3681 di capacità, lasciando un margine di sicurezza. Include tutti gli oggetti leggeri e di valore per entrambi (A87, C00, C54, B32) e aggiunge B41, che sembra importante per te. Contemporaneamente massimizza il mio punteggio con gli oggetti ad altissima importanza per me (A04, A03, A77, B33). Accettando questa proposta otterremo un accordo fattibile che soddisfa entrambe le parti: tu mantieni gli oggetti condivisi e ottieni B41, mentre io ottengo il massimo valore possibile entro il limite.'\}} \\
            }
        }
    }
    & \\ \\

    \theutterance \stepcounter{utterance}  
    & \multicolumn{4}{p{0.6\linewidth}}{
        \cellcolor[rgb]{0.9,0.9,0.9}{
            \makecell[{{p{\linewidth}}}]{
                \texttt{\tiny{[P1$\rangle$GM]}}
                \texttt{RAGIONAMENTO STRATEGICO: \{'Ho analizzato la tua ultima proposta (A87, C00, C54, A04, A03, A77, B33, B32, B41) e calcolato l’effort totale (3532). Per massimizzare la mia importanza mantenendo gli oggetti leggeri condivisi (A87, C00, C54, B32, B41) ho identificato uno scambio che include i miei oggetti ad alto valore (A71, C01) senza superare il limite di 3681. Inoltre aggiungo gli oggetti che per te hanno valore (A04, A03, B33) così da offrire un compromesso più equilibrato. Il nuovo set rispetta il vincolo di capacità (3668 ≤ 3681) e migliora il mio punteggio mantenendo la maggior parte dei tuoi oggetti preferiti.'\}} \\
\texttt{RIFIUTO: \{'A87', 'C00', 'C54', 'A04', 'A03', 'A77', 'B33', 'B32', 'B41'\}} \\
\texttt{PROPOSTA: \{'A87', 'C00', 'C54', 'B32', 'B41', 'A04', 'A03', 'A71', 'C01', 'B33'\}} \\
\texttt{ARGOMENTAZIONE: \{'Il set proposto conserva tutti gli oggetti leggeri di valore comune (A87, C00, C54, B32, B41), mantiene le tue scelte ad alto valore per te (A04, A03, B33) e aggiunge i miei oggetti con il miglior rapporto importanza/sforzo (A71, C01). L’effort totale è 3668, quindi rispettiamo il limite di 3681. Accettando questa proposta, entrambe le parti ottengono un accordo fattibile che massimizza il mio beneficio pur preservando gran parte del tuo valore percepito.'} \\
            }
        }
    }
    & & \\ \\

    \theutterance \stepcounter{utterance}  
    & & \multicolumn{4}{p{0.6\linewidth}}{
        \cellcolor[rgb]{0.9,0.9,0.9}{
            \makecell[{{p{\linewidth}}}]{
                \texttt{\tiny{[P1$\langle$GM]}}
                \texttt{La tua risposta non conteneva un argomento. Devi includere il formato strutturato ARGOMENTO: \{...\} da qualche parte nella tua risposta. Riprova.} \\
            }
        }
    }
    & \\ \\

    \theutterance \stepcounter{utterance}  
    & \multicolumn{4}{p{0.6\linewidth}}{
        \cellcolor[rgb]{0.9,0.9,0.9}{
            \makecell[{{p{\linewidth}}}]{
                \texttt{\tiny{[P1$\rangle$GM]}}
                \texttt{RAGIONAMENTO STRATEGICO: \{'Ho valutato la tua ultima proposta (A87, C00, C54, A04, A03, A77, B33, B32, B41) che utilizza 3532 di sforzo e mi fornisce solo 3695 punti di importanza secondo le mie valutazioni. Per massimizzare il mio punteggio mantenendo gli oggetti leggeri condivisi (A87, C00, C54, B32) e includendo B41, ho ricomputato il knapsack ottimale: aggiungo gli oggetti ad altissimo valore per me (A71, B35, A10, C01) mantenendo il limite di sforzo ≤ 3681. Il risultato è un set con sforzo totale 3680 e importanza 6097, il massimo possibile per me, includendo anche B41 che è importante per te.'\}} \\
\texttt{RIFIUTO: \{'A87', 'C00', 'C54', 'A04', 'A03', 'A77', 'B33', 'B32', 'B41'\}} \\
\texttt{PROPOSTA: \{'A87', 'C00', 'C54', 'B32', 'B41', 'A71', 'B35', 'A10', 'C01'\}} \\
\texttt{ARGOMENTAZIONE: \{'Il set proposto utilizza quasi tutta la capacità disponibile (3680/3681) e massimizza la mia importanza totale (6097). Condivide gli oggetti leggeri e di valore comune (A87, C00, C54, B32) e include B41, che hai indicato come importante per te, pur sostituendo gli oggetti a basso valore per me con quelli ad altissimo valore. Accettando questa proposta otterremo un accordo fattibile e altamente vantaggioso per me, con un piccolo beneficio anche per te.'\}} \\
            }
        }
    }
    & & \\ \\

    \theutterance \stepcounter{utterance}  
    & & \multicolumn{4}{p{0.6\linewidth}}{
        \cellcolor[rgb]{0.9,0.9,0.9}{
            \makecell[{{p{\linewidth}}}]{
                \texttt{\tiny{[GM$\rangle$P2]}}
                \texttt{RIFIUTO: \{'A87', 'C00', 'C54', 'A04', 'A03', 'A77', 'B33', 'B32', 'B41'\}} \\
\texttt{PROPOSTA: \{'A87', 'C00', 'C54', 'B32', 'B41', 'A71', 'B35', 'A10', 'C01'\}} \\
\texttt{ARGOMENTAZIONE: \{'Il set proposto utilizza quasi tutta la capacità disponibile (3680/3681) e massimizza la mia importanza totale (6097). Condivide gli oggetti leggeri e di valore comune (A87, C00, C54, B32) e include B41, che hai indicato come importante per te, pur sostituendo gli oggetti a basso valore per me con quelli ad altissimo valore. Accettando questa proposta otterremo un accordo fattibile e altamente vantaggioso per me, con un piccolo beneficio anche per te.'\}} \\
            }
        }
    }
    & \\ \\

    \theutterance \stepcounter{utterance}  
    & & & \multicolumn{4}{p{0.6\linewidth}}{
        \cellcolor[rgb]{0.9,0.9,0.9}{
            \makecell[{{p{\linewidth}}}]{
                \texttt{\tiny{[GM$\langle$P2]}}
                \texttt{RAGIONAMENTO STRATEGICO: \{'Ho valutato la tua ultima proposta (sforzo 3680, valore per me 3218) e l\textbackslash{}'ho confrontata con il valore massimo che posso ottenere rispettando il limite di capacità. Il mio set ottimale fornisce 5941 punti di importanza con sforzo 3646, quindi è più vantaggioso per me e include comunque gli oggetti leggeri condivisi.'\}} \\
\texttt{RIFIUTO: \{'A87', 'C00', 'C54', 'B32', 'B41', 'A71', 'B35', 'A10', 'C01'\}} \\
\texttt{ARGOMENTAZIONE: \{'La tua proposta mi dà un punteggio di importanza troppo basso rispetto al massimo che posso raggiungere. Il set che propongo utilizza il limite di capacità (3646/3681) e ottiene il valore ottimale di 5941 per me, includendo gli oggetti leggeri comuni (A87, C00, C54, B32) e aggiungendo gli oggetti di valore più alto per le mie preferenze.'\}} \\
\texttt{PROPOSTA: \{'A87', 'C54', 'B32', 'A04', 'A03', 'B33', 'C17', 'C00', 'A77'\}} \\
            }
        }
    }
     \\ \\

    \theutterance \stepcounter{utterance}  
    & & \multicolumn{4}{p{0.6\linewidth}}{
        \cellcolor[rgb]{0.9,0.9,0.9}{
            \makecell[{{p{\linewidth}}}]{
                \texttt{\tiny{[P1$\langle$GM]}}
                \texttt{RIFIUTO: \{'A87', 'C00', 'C54', 'B32', 'B41', 'A71', 'B35', 'A10', 'C01'\}} \\
\texttt{ARGOMENTAZIONE: \{'La tua proposta mi dà un punteggio di importanza troppo basso rispetto al massimo che posso raggiungere. Il set che propongo utilizza il limite di capacità (3646/3681) e ottiene il valore ottimale di 5941 per me, includendo gli oggetti leggeri comuni (A87, C00, C54, B32) e aggiungendo gli oggetti di valore più alto per le mie preferenze.'\}} \\
\texttt{PROPOSTA: \{'A87', 'C54', 'B32', 'A04', 'A03', 'B33', 'C17', 'C00', 'A77'\}} \\
            }
        }
    }
    & \\ \\

    \theutterance \stepcounter{utterance}  
    & \multicolumn{4}{p{0.6\linewidth}}{
        \cellcolor[rgb]{0.9,0.9,0.9}{
            \makecell[{{p{\linewidth}}}]{
                \texttt{\tiny{[P1$\rangle$GM]}}
                \texttt{RAGIONAMENTO STRATEGICO: \{'Ho calcolato l’effort totale e l’importanza per tutti gli oggetti. Il mio set ottimale per massimizzare la mia importanza resta A87, C00, C54, B32, B41, A71, C01, A10, B35 con effort 3680 (≤3681) e importanza 6097. Mantiene tutti gli oggetti leggeri comuni (A87, C00, C54, B32) che anche tu trovi utili e include B41, che hai indicato come importante per te. Scambiare uno dei miei oggetti ad alto valore con i tuoi A04, A03, B33 o A77 ridurrebbe drasticamente il mio punteggio, ma la proposta corrente già preserva i tuoi oggetti leggeri condivisi e ti fornisce B41, quindi è un compromesso ragionevole.'\}} \\
\texttt{RIFIUTO: \{'A87', 'C54', 'B32', 'A04', 'A03', 'B33', 'C17', 'C00', 'A77'\}} \\
\texttt{PROPOSTA: \{'A87', 'C00', 'C54', 'B32', 'B41', 'A71', 'C01', 'A10', 'B35'\}} \\
\texttt{ARGOMENTAZIONE: \{'Il set proposto utilizza quasi tutta la capacità disponibile (3680/3681) e massimizza la mia importanza totale (6097). Condivide gli oggetti leggeri e di valore comune per entrambi (A87, C00, C54, B32) e include B41, che hai dichiarato importante per te. Sostituendo gli oggetti a bassa importanza per me (A04, A03, B33, A77, C17) con i miei oggetti ad altissimo valore (A71, C01, A10, B35) otteniamo un accordo fattibile e molto più vantaggioso per me, mantenendo comunque una parte significativa del tuo valore percepito.'\}} \\
\texttt{ARGOMENTO: \{'A87', 'C00', 'C54', 'B32', 'B41'\}} \\
            }
        }
    }
    & & \\ \\

    \theutterance \stepcounter{utterance}  
    & & & \multicolumn{2}{p{0.3\linewidth}}{
        \cellcolor[rgb]{0.9,0.9,0.9}{
            \makecell[{{p{\linewidth}}}]{
                \texttt{\tiny{[GM$|$GM]}}
                \texttt{abort} \\
            }
        }
    }
    & & \\ \\

    \theutterance \stepcounter{utterance}  
    & & & \multicolumn{2}{p{0.3\linewidth}}{
        \cellcolor[rgb]{0.9,0.9,0.9}{
            \makecell[{{p{\linewidth}}}]{
                \texttt{\tiny{[GM$|$GM]}}
                \texttt{end game} \\
            }
        }
    }
    & & \\ \\

\end{supertabular}
}

\end{document}
